% Chapter 3: Overall Structure
\chapter{总体架构与章节体系对比}

\section{章节结构总览}

\subsection{章节数量对比}

\begin{table}[H]
\centering
\caption{三版本章节数量对比}
\begin{tabular}{|l|c|c|c|}
\hline
\textbf{项目} & \textbf{\versionAShort} & \textbf{\versionBShort} & \textbf{\versionCShort} \\
\hline
总章节数     & 12 章 + 8 附录     & 13 章 + 9 附录     & 13 章 + 11 附录     \\
\hline
正文页数     & 约 90 页           & 约 110 页          & 约 130 页          \\
\hline
附录页数     & 约 60 页           & 约 80 页           & 约 120 页          \\
\hline
总页数       & 约 150 页          & 约 190 页          & 约 250 页          \\
\hline
表格数       & 约 50 个           & 约 70 个           & 约 90 个           \\
\hline
图示数       & 约 30 个           & 约 40 个           & 约 60 个           \\
\hline
\end{tabular}
\end{table}

\subsection{章节标题对照表}

下表给出三版本主要章节的对照（基于实际法规文本）：

\begin{table}[H]
\centering
\caption{三版本主要章节对照（基于实际法规文本）}
\small
\begin{tabularx}{\textwidth}{|c|X|X|X|}
\hline
\textbf{序号} & \textbf{\versionAShort} & \textbf{\versionBShort} & \textbf{\versionCShort} \\
\hline
1 & 范围                   & 范围                   & Scope \\
\hline
2 & 规范性引用文件         & 规范性引用文件         & Definitions \\
\hline
3 & 术语和定义             & 术语和定义             & Application for approval \\
\hline
4 & 一般规定               & 一般规定               & Approval \\
\hline
5 & 特殊规定               & 主功能的特殊规定       & General specifications \\
\hline
  & 5.1 远光灯             & 5 远光灯（无自适应模式）  & 6 Individual specifications \\
\hline
  & 5.2 近光灯             & 5 近光灯（无自适应模式）  & \\
\hline
  & 5.3 前雾灯             & 5 带有自适应模式的前照灯（近光灯和/或远光灯） & \\
\hline
  & 5.4 倒车灯             & 5 前雾灯 & \\
\hline
  & 5.5 转向信号灯         & 5 角灯 & \\
\hline
  & 5.6 危险警告信号       & 5 转向信号灯（转向灯）   & \\
\hline
  & 5.7 制动灯             & 5 制动灯             & \\
\hline
  & 5.8 后牌照板照明装置   & 5 后牌照板照明装置      & \\
\hline
  & 5.9 前位灯             & 5 前位灯             & \\
\hline
  & 5.10 后位灯            & 5 后位灯             & \\
\hline
  & 5.11 后雾灯            & 5 后雾灯             & \\
\hline
  & 5.12 驻车灯            & 5 驻车灯             & \\
\hline
  & 5.13 示廓灯            & 5 示廓灯             & \\
\hline
  & 5.14/15/16/17 反射器   & 5 非三角形/三角形后回复反射器 & \\
\hline
  & 5.18 侧标志灯          & 5 非三角形前/侧回复反射器 & \\
\hline
  & 5.19 昼间行驶灯        & 5 侧标志灯             & \\
\hline
  & 5.20 角灯              & 5 昼间行驶灯           & \\
\hline
  & 5.21 醒目标志          & 5 醒目标志             & \\
\hline
  & 5.22 自适应前照明系统(AFS) & 5 倒车灯          & \\
\hline
  & 5.23 紧急制动信号      & 5 慢行灯（新增）        & \\
\hline
  & 5.24 牵引杆挂车标志灯  & 5 牵引杆挂车标志灯      & \\
\hline
  & 5.25 慢行灯            & \textcolor{themegreen}{5 自动化驾驶系统标志灯（ADS 标志灯）（新增）} & \\
\hline
  & 5.26 追尾警告信号      & 6 从属功能的特殊规定（新增） & \\
\hline
  &                        & 6 驾驶员辅助投射功能     & \\
\hline
  &                        & 6 光信号投射功能         & \\
\hline
  &                        & 6 图文交互               & \\
\hline
  &                        & 6 发光徽标               & \\
\hline
6 & 检验规则               & 7 扩展功能的特殊规定（新增） & 7 Modifications and extensions of approval \\
\hline
  &                        & 7 危险警告信号           & 8 Conformity of production \\
\hline
  &                        & 7 紧急制动信号           & 9 Penalties for non-conformity \\
\hline
  &                        & 7 追尾警告信号           & 10 Production definitively discontinued \\
\hline
  &                        & 8 检验规则               & 11 Names and addresses \\
\hline
  &                        & \textcolor{themegreen}{9 同一型式判定（新增）} & 12 Transitional provisions \\
\hline
  &                        & \textcolor{themegreen}{10 实施日期（新增）}  & \\
\hline
\end{tabularx}
\end{table}

\section{草案版的体系调整}

\subsection{新增章节}

草案版相对现行版新增了以下章节：

\begin{enumerate}[leftmargin=2em]
    \item \textbf{第 4 章 新增小节 4.7 "电磁兼容性（EMC）"}：明确灯具的电磁兼容要求；
    \item \textbf{第 5 章 新增 5.4 "ADB 自适应远光灯"}：详细规定 ADB 的安装、控制、测试；
    \item \textbf{第 5 章 新增 5.5 "矩阵式 LED 前照灯"}：针对矩阵式 LED 的特殊规定；
    \item \textbf{第 7 章 新增 7.6 "Logo 投影灯与迎宾灯"}：对新型装饰性灯具作出规定；
    \item \textbf{新增第 12 章 "智能照明系统"}：综合论述 ADB、AFS、Logo 灯、迎宾灯等智能化照明；
    \item \textbf{新增附录 D "ADB 系统测试方法"}：详细的 ADB 测试程序；
    \item \textbf{新增附录 E "智能照明故障诊断"}：智能照明系统的故障诊断与指示要求。
\end{enumerate}

\subsection{调整章节}

草案版对以下章节进行了调整：

\begin{enumerate}[leftmargin=2em]
    \item \textbf{第 4 章 一般规定}：拆分为"基本要求"（4.1--4.5）和"特殊要求"（4.6--4.10）；
    \item \textbf{第 6 章 信号灯装置}：内部小节重新编排，按"前部信号灯"、"后部信号灯"、"侧部信号灯"重新组织；
    \item \textbf{第 10 章 特殊装置}：将 ESS 升级为独立小节，并增加"接近警告灯"等新型装置；
    \item \textbf{第 11 章 认证与标志}：增加电子标识、远程认证等内容。
\end{enumerate}

\subsection{删除章节}

草案版删除了现行版中的部分章节（内容整合到其他章节）：

\begin{itemize}[leftmargin=2em]
    \item 原"测试方法"独立章节，并入相关具体装置章节；
    \item 部分重复性附录合并；
    \item 老旧测试设备相关附录删除。
\end{itemize}

\section{UN R48-09 的体系特点}

\subsection{逻辑层次}

UN R48-09 在结构上采用"通用要求 $\rightarrow$ 通用装置 $\rightarrow$ 特定装置 $\rightarrow$ 测试方法 $\rightarrow$ 认证程序"的逻辑层次：

\begin{enumerate}[leftmargin=2em]
    \item \textbf{基础部分}：Scope、Definitions、Application of regulations；
    \item \textbf{通用规格}：General specifications（适用于所有装置）；
    \item \textbf{通用装置规格}：Headlamps、Direction indicators 等；
    \item \textbf{特殊装置规格}：Daytime running lamps、Adaptive front-lighting 等；
    \item \textbf{认证要求}：Approval and installation requirements；
    \item \textbf{附录}：详细测试程序、数据表、示意图。
\end{enumerate}

\subsection{编排优势}

UN R48-09 的体系安排具有以下优势：

\begin{itemize}[leftmargin=2em]
    \item \textbf{层次清晰}：从一般到特殊的递进式结构；
    \item \textbf{模块化}：各装置独立成章，便于修订和增补；
    \item \textbf{索引便利}：详尽的章节索引和交叉引用；
    \item \textbf{附录丰富}：附录包含大量技术细节，便于查询。
\end{itemize}

\subsection{对 GB 4785 草案的影响}

UN R48-09 的体系对 GB 4785 草案产生了明显影响：

\begin{enumerate}[leftmargin=2em]
    \item 草案版引入了"通用规格"章节（与 UN R48-09 的第 4 章对应）；
    \item 草案版将信号灯装置按"前/后/侧"重新组织，借鉴了 UN R48 的做法；
    \item 草案版的附录数量和内容向 UN R48-09 看齐；
    \item 草案版对术语章节进行了扩充，增加了大量新术语。
\end{enumerate}

\section{编排逻辑差异分析}

\subsection{中国特色编排}

GB 4785 的编排体系在某些方面体现了中国特色：

\begin{enumerate}[leftmargin=2em]
    \item \textbf{强调基础通用性}：将"基本要求"和"具体要求"分章处理，更便于产品认证；
    \item \textbf{重视标志标识}：将标志标识作为独立章节，便于管理；
    \item \textbf{附录实用化}：附录更注重表格化、操作性强的内容。
\end{enumerate}

\subsection{欧洲编排特点}

UN R48-09 的编排更注重：

\begin{enumerate}[leftmargin=2em]
    \item \textbf{法规的逻辑完整性}：从定义到认证形成完整闭环；
    \item \textbf{技术细节的可追溯性}：通过交叉引用和注释保证可追溯；
    \item \textbf{对新技术的高度开放性}：预留扩展空间，便于未来修订。
\end{enumerate}

\section{章节对照矩阵}

\subsection{主要章节对照矩阵}

\begin{table}[H]
\centering
\caption{主要章节对照矩阵}
\small
\begin{tabularx}{\textwidth}{|l|c|c|c|c|}
\hline
\textbf{主题} & \textbf{\versionAShort} & \textbf{\versionBShort} & \textbf{\versionCShort} & \textbf{对齐度} \\
\hline
适用范围         & 1    & 1    & 1     & 高 \\
\hline
规范性引用文件   & 2    & 2    & ---   & --- \\
\hline
术语和定义       & 3    & 3    & 2     & 高 \\
\hline
一般规定         & 4    & 4    & 4     & 高 \\
\hline
前照灯（远/近）  & 5.1--5.2 & 5.1--5.3 & 5.1--5.6 & 高 \\
\hline
AFS              & 5.3  & 5.4  & 6     & 中 \\
\hline
              & ---  & 5.5  & 6     & 高（草案与欧版） \\
\hline
信号灯           & 6    & 6    & 6     & 高 \\
\hline
位置/示廓灯     & 7.1--7.2 & 7.1--7.3 & 6     & 高 \\
\hline
标志与反射器     & 8    & 8    & 7     & 高 \\
\hline
牌照灯          & 9    & 9    & 6     & 中 \\
\hline
雾灯（前/后）   & 10   & 10   & 5     & 中 \\
\hline
特殊装置        & 11   & 11    & 11    & 中 \\
\hline
ESS             & 11   & 11    & 11    & 高 \\
\hline
认证            & 12   & 12    & 12    & 高 \\
\hline
测试方法        & 附录  & 附录  & 附录  & 高 \\
\hline
\end{tabularx}
\end{table}

\subsection{新增/删除章节对照}

\begin{table}[H]
\centering
\caption{草案版新增/删除章节}
\small
\begin{tabularx}{\textwidth}{|l|X|c|}
\hline
\textbf{变更} & \textbf{内容描述} & \textbf{对标欧版} \\
\hline
新增 5.4 ADB 系统       & 新增自适应远光灯系统独立规定      & R48 §6 \\
\hline
新增 5.5 矩阵式 LED      & 新增矩阵式 LED 前照灯规定        & R48 §6 \\
\hline
新增 7.6 Logo 灯/迎宾灯   & 新增装饰性灯具规定              & R48 §6 \\
\hline
新增 12 智能照明系统     & 新增智能照明综合章节            & R48 §12 \\
\hline
新增 附录 D ADB 测试     & 新增 ADB 测试程序              & R48 Annex \\
\hline
新增 附录 E 智能诊断     & 新增智能诊断要求                & R48 Annex \\
\hline
删除 原"测试方法"独立章  & 合并入相关装置章节            & --- \\
\hline
\end{tabularx}
\end{table}

\section{小结}

\begin{enumerate}[leftmargin=2em]
    \item 草案版在总体框架上与现行版保持高度一致，但通过新增章节适应新技术的发展；
    \item 草案版在编排逻辑上明显借鉴了 UN R48-09，特别是"前/后/侧"的分类方式；
    \item 草案版新增的 7 个章节（或小节）均对应 UN R48-09 中的成熟条款，反映了国际化对接的明确方向；
    \item 草案版相对现行版的体系调整体现了"以我为主、借鉴国际"的中国标准化工作思路。
\end{enumerate}